%-------------------------
% 夏驭风 —— 中文简历（AI infra 与芯片架构工程师 · 超芯人才计划）
% 共享导言区，由 yufeng_cn_1p.tex / yufeng_cn_2p.tex 引入
%-------------------------

\usepackage[UTF8]{ctex}
\usepackage[empty]{fullpage}
\usepackage{titlesec}
\usepackage[usenames,dvipsnames]{color}
\usepackage{enumitem}
\usepackage[hidelinks]{hyperref}
\usepackage{fancyhdr}
\usepackage{tabularx}

\pagestyle{fancy}
\fancyhf{}
\fancyfoot{}
\renewcommand{\headrulewidth}{0pt}
\renewcommand{\footrulewidth}{0pt}

% 页边距
\addtolength{\oddsidemargin}{-0.6in}
\addtolength{\evensidemargin}{-0.6in}
\addtolength{\textwidth}{1.2in}
\addtolength{\topmargin}{-.8in}
\addtolength{\textheight}{1.6in}

% 列表间距收紧
\setlist[itemize]{itemsep=0pt,parsep=0pt,topsep=2pt,partopsep=0pt}

\urlstyle{same}
\raggedbottom
\raggedright
\setlength{\tabcolsep}{0in}

% 章节格式
\titleformat{\section}{
  \scshape\raggedright\normalsize\bfseries
}{}{0em}{}[\color{black}\titlerule]
\titlespacing*{\section}{0pt}{6pt}{3pt}

%-------------------------
% 自定义命令
\newcommand{\resumeItem}[2]{
  \item\small{
    \textbf{#1}{：#2 \vspace{-2pt}}
  }
}

\newcommand{\resumeItemPlain}[1]{
  \item\small{{#1 \vspace{-2pt}}}
}

% 一行式条目：标签 + 内容（技能栏、荣誉栏）
\newcommand{\skillrow}[2]{
  \item\small{\textbf{#1}\hspace{0.4em}#2 \vspace{-4pt}}
}

\newcommand{\resumeSubheading}[4]{
  \vspace{-1pt}\item
    \begin{tabular*}{0.97\textwidth}[t]{@{}p{0.72\textwidth}@{\extracolsep{\fill}}r@{}}
      \textbf{#1} & #2 \\
      \textit{\small #3} & \textit{\small #4} \\
    \end{tabular*}\vspace{-5pt}
}

% 紧凑条目：标题 + 右侧信息，正文另起一段（用于无子项的课程项目）
\newcommand{\resumeCompact}[3]{
  \vspace{-1pt}\item
    \begin{tabular*}{0.97\textwidth}[t]{@{}p{0.72\textwidth}@{\extracolsep{\fill}}r@{}}
      \textbf{#1} & #2 \\
    \end{tabular*}\par\vspace{1pt}
    {\small #3}\vspace{-1pt}
}

% 论文条目：标题（左，加粗）+ 会议（右），作者/位次另起一行
\newcommand{\pubentry}[3]{%
  \begin{tabular*}{\textwidth}[t]{@{}p{0.76\textwidth}@{\extracolsep{\fill}}r@{}}
    \small\textbf{#1} & \small\textit{#2}
  \end{tabular*}\par\vspace{-2pt}
  {\small\textit{#3}}\par\vspace{3pt}
}

\newcommand{\resumeSubItem}[2]{\resumeItem{#1}{#2}\vspace{-4pt}}

\renewcommand{\labelitemii}{$\circ$}

\newcommand{\resumeSubHeadingListStart}{\begin{itemize}[leftmargin=*]}
\newcommand{\resumeSubHeadingListEnd}{\end{itemize}}
\newcommand{\resumeItemListStart}{\begin{itemize}}
\newcommand{\resumeItemListEnd}{\end{itemize}\vspace{-5pt}}
